% ============================================================
% TOUCH — Section 3.0: Haptic Epistemology
% ============================================================
% Laura Marks — The Skin of the Film: haptic visuality;
% the skin as an epistemological organ; touch as a way of
% knowing that exceeds and precedes vision.
% The HairSynth is a haptic epistemology instrument: it
% makes tactile knowledge into audible form.
% The data-sound mappings: not "more texture = louder."
% The mappings encode embodied knowledge — what it feels
% like to touch hair. Emerged through embodied
% experimentation, adjusting while listening until the
% output "felt right." This is methodology.
% Agential Realism (Barad): the HairSynth as apparatus.
% Hair data + synthesis system intra-act to produce the
% phenomenon (sound of hair) that neither produces alone.
% Sonic virtuality (Grimshaw & Garner): the drone returns
% explicitly here. Sound was the condition of possibility
% before any of this was built.
%
% Key texts:
%   Marks — The Skin of the Film: Intercultural Cinema,
%   Embodiment, and the Senses
%   Barad — Meeting the Universe Halfway (returns here)
%   Grimshaw & Garner — Sonic Virtuality (drone explicit)
% ============================================================

\section*{Section 3.0: Haptic Epistemology}

\ifdraft
\lipsum[1-3]
\fi
